% !PRINTED 47
\weberpage{47}

Wir wollen das gewonnene Resultat noch als Satz aussprechen:

5. \gesperrt{Die Charaktere einer Gruppe $S$ können zu einer mit $S$ isomorphen Gruppe verbunden werden.}

Es sei noch der Satz erwähnt, der sich aus der Definition von $\chi_B$ \gesperrt{durch} die Formel (13) unmittelbar ergiebt:
\weq{20}{\chi_B(A) = \chi_A(B).}

Endlich gilt auch noch der folgende Satz:

6. \gesperrt{Durchläuft $A$ alle Elemente der Gruppe $S$, so ist für einen feststehenden Charakter $\chi$:}
\weq{21}{\overset{A}{\Sigma}\,\chi(A) = n \text{ oder } = 0,}
\gesperrt{je nachdem $\chi$ der Einheitscharakter ist oder nicht. Ebenso ist, wenn das Element $A$ festgehalten wird und $\chi$ die Reihe der Charaktere durchläuft:}
\weq{22}{\overset{\chi}{\Sigma}\,\chi(A) = n \text{ oder } = 0,}
\gesperrt{je nachdem $A$ das Einheitselement ist oder nicht.}

Für den Fall, dass $\chi$ oder $A$ die Einheitselemente ihrer Gruppen sind, sind die Formeln (21) und (22) evident, da dann jedes der $n$ Glieder der Summe den Werth 1 hat. Ist aber $\chi$ nicht der Einheitscharakter, so giebt es ein Element $B$ in $S$, so dass $\chi(B)$ nicht $= 1$ ist. Setzen wir dann
\[
\overset{A}{\Sigma}\,\chi(A) = \sigma,
\]
so folgt durch Multiplication mit $\chi(B)$:
\[
\overset{A}{\Sigma}\,\chi(AB) = \sigma\chi(B).
\]
Da aber $AB$ zugleich mit $A$ die ganze Gruppe $S$ durchläuft, so ist auch $\overset{A}{\Sigma}\,\chi(AB) = \sigma$, also $\sigma[1 - \chi(B)] = 0$ oder $\sigma = 0$, w.\,z.\,b.\,w.

Ebenso beweist man die Formel (22), die übrigens auch nach (20) unmittelbar aus (21) folgt.
